% Emacs settings: -*-mode: latex; TeX-master: "Component_manual"; -*-

\chapter{The Union framework: sample environments and multiple scattering}
\label{c:union-comp}
\index{Library!Components!union}
\index{Union|textbf}

The Union framework, originally contributed by Mads Bertelsen for \MCS
and ported to \MCX, is a set of components that together allow the
description of complex, possibly nested or overlapping, sample
environments with multiple scattering, by separating the description of
\emph{geometry} (where matter is) from \emph{physical process} (what
happens to a photon there, e.g. Compton or Rayleigh scattering,
photoelectric absorption, or powder diffraction). For a conceptual
introduction to the framework -- volumes, priority, masks, exit volumes,
and the underlying propagation algorithm -- see chapter~\ref{c:union}.
This chapter documents the input parameters of every individual Union
component.

\section{Union framework core: initialization, master, and termination}
\input{union/Union_init}
\input{union/Union_master}
\section{The \texttt{Union\_stop} McXtrace Component}
Stop component that must be placed after all Union components for instrument to compile correctly

\subsection*{Identification}
\begin{itemize}
  \item \textbf{Author:} Mads Bertelsen and Erik B Knudsen
  \item \textbf{Origin:} ESS DMSC \& DTU Physics
  \item \textbf{Date:} 20.08.15
\end{itemize}

\subsection*{Description}
Part of the Union components, a set of components that work together and thus separates geometry and physics within McXtrace. The use of this component requires other components to be used.

1) One specifies a number of processes using process components 2) These are gathered into material definitions using Union\_make\_material 3) Geometries are placed using Union\_box/cylinder/sphere, assigned a material 4) A Union\_master component placed after all of the above

Only in step 4 will any simulation happen, and per default all geometries defined before this master, but after the previous will be simulated here.

There is a dedicated manual available for the Union components

Algorithm: Described elsewhere

\subsection*{Input parameters}
Parameters in \textbf{boldface} are required; the others are optional.

\begin{longtable}{p{0.22\textwidth}p{0.12\textwidth}p{0.46\textwidth}p{0.14\textwidth}}
\toprule
\textbf{Name} & \textbf{Unit} & \textbf{Description} & \textbf{Default} \\
\midrule
\endhead
\bottomrule
\end{longtable}

\subsection*{Links}
\begin{itemize}
  \item Component source code found in file \texttt{Union\_stop.comp}.
\end{itemize}
\IfFileExists{union/Union_stop_static.tex}{\input{union/Union_stop_static.tex}}{}

\section{Union geometry components}
\section{The \texttt{Union\_box} McXtrace Component}
Implementation of a box geometry - to be filled with a material

\subsection*{Identification}
\begin{itemize}
  \item \textbf{Author:} Mads Bertelsen and Erik B Knudsen
  \item \textbf{Origin:} ESS DMSC \& DTU Physics
  \item \textbf{Date:} 20.08.15
\end{itemize}

\subsection*{Description}
Part of the Union components, a set of components that work together and thus separates geometry and physics within McXtrace. The use of this component requires other components to be used.

1) One specifies a number of processes using process components 2) These are gathered into material definitions using Union\_make\_material 3) Geometries are placed using Union\_box/cylinder/sphere, assigned a material 4) A Union\_master component placed after all of the above

Only in step 4 will any simulation happen, and per default all geometries defined before this master, but after the previous will be simulated here.

There is a dedicated manual available for the Union components

The position of this component is the center of the box, extending xwidth/2, yheight/2, and zdepth/2 in each direction respectively.

It is allowed to overlap components, but it is not allowed to have two parallel planes that coincide. This will crash the code on run time.

Algorithm: Described elsewhere

\subsection*{Input parameters}
Parameters in \textbf{boldface} are required; the others are optional.

\begin{longtable}{p{0.22\textwidth}p{0.12\textwidth}p{0.46\textwidth}p{0.14\textwidth}}
\toprule
\textbf{Name} & \textbf{Unit} & \textbf{Description} & \textbf{Default} \\
\midrule
\endhead
material\_string & str & material name of this volume, defined using Union\_make\_material & 0 \\
\textbf{priority} & 1 & priotiry of the volume (can not be the same as another volume) A high priority is on top of low. &  \\
\textbf{xwidth} & m & width of the box volume &  \\
\textbf{yheight} & m & height of the box volume &  \\
\textbf{zdepth} & m & depth of the box volume &  \\
xwidth2 & m & optional different width at the +z box face & -1 \\
yheight2 & m & optional different height at the +z box face & -1 \\
visualize & 1 & set to 0 if you wish to hide this geometry in mcdisplay & 1 \\
target\_index & 1 & Focuses on component a component this many steps further in the component sequence & 0 \\
target\_x & m & \textbackslash{} & 0 \\
target\_y & m & - Position of target to focus at & 0 \\
target\_z & m & / & 0 \\
focus\_aw & deg & horiz. angular dimension of a rectangular area & 0 \\
focus\_ah & deg & vert. angular dimension of a rectangular area & 0 \\
focus\_xw & m & horiz. dimension of a rectangular area & 0 \\
focus\_xh & m & vert. dimension of a rectangular area & 0 \\
focus\_r & m & focusing on circle with this radius & 0 \\
p\_interact & 1 & probability to interact with this geometry [0-1] & 0 \\
mask\_string & str & Comma separated list of geometry names which this geometry should mask & 0 \\
mask\_setting & str & "All" or "Any", should the masked volume be simulated when the ray is in just one mask, or all. & 0 \\
number\_of\_activations & 1 & Number of subsequent Union\_master components that will simulate this geometry & 1 \\
init & str & Name of Union inititaliser component & "init" \\
\bottomrule
\end{longtable}

\subsection*{Links}
\begin{itemize}
  \item Component source code found in file \texttt{Union\_box.comp}.
\end{itemize}
\IfFileExists{union/Union_box_static.tex}{\input{union/Union_box_static.tex}}{}
\input{union/Union_cylinder}
\input{union/Union_sphere}
\section{The \texttt{Union\_cone} McXtrace Component}
Cone geometry component for Union components

\subsection*{Identification}
\begin{itemize}
  \item \textbf{Author:} Mads Bertelsen and Erik B Knudsen
  \item \textbf{Origin:} ESS DMSC \& DTU Physics
  \item \textbf{Date:} 20.08.15
\end{itemize}

\subsection*{Description}
Part of the Union components, a set of components that work together and thus separates geometry and physics within McXtrace. The use of this component requires other components to be used.

1) One specifies a number of processes using process components 2) These are gathered into material definitions using Union\_make\_material 3) Geometries are placed using Union\_box/cylinder/sphere, assigned a material 4) A Union\_master component placed after all of the above

Only in step 4 will any simulation happen, and per default all geometries defined before this master, but after the previous will be simulated here.

There is a dedicated manual available for the Union components

The position of this component is the center of the cone, and it thus extends yheight/2 up and down along y axis.

It is allowed to overlap components, but it is not allowed to have two parallel planes that coincide. This will crash the code on run time.

\subsection*{Input parameters}
Parameters in \textbf{boldface} are required; the others are optional.

\begin{longtable}{p{0.22\textwidth}p{0.12\textwidth}p{0.46\textwidth}p{0.14\textwidth}}
\toprule
\textbf{Name} & \textbf{Unit} & \textbf{Description} & \textbf{Default} \\
\midrule
\endhead
material\_string & string & Material name of this volume, defined using Union\_make\_material & 0 \\
\textbf{priority} & 1 & Priotiry of the volume (can not be the same as another volume) A high priority is on top of low. &  \\
radius & m & Radius volume in (x,z) plane & 0 \\
radius\_top & m & Top radius volume in (x,z) plane & 0 \\
radius\_bottom & m & Bottom radius volume in (x,z) plane & 0 \\
\textbf{yheight} & m & Cone height in (y) direction &  \\
visualize & 1 & Set to 0 if you wish to hide this geometry in mcdisplay & 1 \\
target\_index & 1 & Focuses on component a component this many steps further in the component sequence & 0 \\
target\_x & m & \textbackslash{} & 0 \\
target\_y & m & - Position of target to focus at & 0 \\
target\_z & m & / & 0 \\
focus\_aw & deg & Horiz. angular dimension of a rectangular area & 0 \\
focus\_ah & deg & Vert. angular dimension of a rectangular area & 0 \\
focus\_xw & m & Horiz. dimension of a rectangular area & 0 \\
focus\_xh & m & Vert. dimension of a rectangular area & 0 \\
focus\_r & m & Focusing on circle with this radius & 0 \\
p\_interact & 1 & Probability to interact with this geometry [0-1] & 0 \\
mask\_string & string & Comma seperated list of geometry names which this geometry should mask & 0 \\
mask\_setting & string & "All" or "Any", should the masked volume be simulated when the ray is in just one mask, or all. & 0 \\
number\_of\_activations & 1 & Number of subsequent Union\_master components that will simulate this geometry & 1 \\
init & string & Name of Union\_init component (typically "init", default) & "init" \\
\bottomrule
\end{longtable}

\subsection*{Links}
\begin{itemize}
  \item Component source code found in file \texttt{Union\_cone.comp}.
\end{itemize}
\IfFileExists{union/Union_cone_static.tex}{\input{union/Union_cone_static.tex}}{}

\section{Union material definition}
\section{The \texttt{Union\_make\_material} McXtrace Component}
Component that takes a number of Union processes and constructs a Union material

\subsection*{Identification}
\begin{itemize}
  \item \textbf{Author:} Mads Bertelsen and Erik B Knudsen
  \item \textbf{Origin:} ESS DMSC \& DTU Physics
  \item \textbf{Date:} 20.08.15
\end{itemize}

\subsection*{Description}
Part of the Union components, a set of components that work together and thus separates geometry and physics within McXtrace. The use of this component requires other components to be used.

1) One specifies a number of processes using process components 2) These are gathered into material definitions using this component 3) Geometries are placed using Union\_box/cylinder/sphere, assigned a material 4) A Union\_master component placed after all of the above

Only in step 4 will any simulation happen, and per default all geometries defined before the master, but after the previous will be simulated here.

There is a dedicated manual available for the Union\_components

Algorithm: Described elsewhere

\subsection*{Input parameters}
Parameters in \textbf{boldface} are required; the others are optional.

\begin{longtable}{p{0.22\textwidth}p{0.12\textwidth}p{0.46\textwidth}p{0.14\textwidth}}
\toprule
\textbf{Name} & \textbf{Unit} & \textbf{Description} & \textbf{Default} \\
\midrule
\endhead
process\_string & string & Comma seperated names of physical processes & "NULL" \\
\textbf{my\_absorption} & 1/m & Inverse penetration depth from absorption at standard energy &  \\
absorber & 0/1 & Control parameter, if set to 1 the material will have no scattering processes & 0 \\
material\_string & string & List of elements present in the material. Triggers a search for materials constants files. & "NULL" \\
init & string & Name of Union\_init component (typically "init", default) & "init" \\
\bottomrule
\end{longtable}

\subsection*{Links}
\begin{itemize}
  \item Component source code found in file \texttt{Union\_make\_material.comp}.
\end{itemize}
\IfFileExists{union/Union_make_material_static.tex}{\input{union/Union_make_material_static.tex}}{}

\section{Union scattering/interaction process components}
\input{union/Incoherent_process}
\input{union/Compton_xrl_process}
\input{union/Rayleigh_xrl_process}
\section{The \texttt{KN\_xrl\_process} McXtrace Component}
A component implementing the Klein-Nishina cross section as a Union physics process

\subsection*{Identification}
\begin{itemize}
  \item \textbf{Author:} Mads Bertelsen and Erik B Knudsen
  \item \textbf{Origin:} ESS DMSC \& DTU Physics \& United Neux
  \item \textbf{Date:} 20.08.15
\end{itemize}

\subsection*{Description}
This Union\_process is based on the Incoherent.comp component originally written by Kim Lefmann and Kristian Nielsen

Part of the Union components, a set of components that work together and thus separates geometry and physics within McXtrace. The use of this component requires other components to be used.

1) One specifies a number of processes using process components like this one 2) These are gathered into material definitions using Union\_make\_material 3) Geometries are placed using Union\_box / Union\_cylinder, assigned a material 4) A Union\_master component placed after all of the above

Only in step 4 will any simulation happen, and per default all geometries defined before the master, but after the previous will be simulated here.

There is a dedicated manual available for the Union\_components

Algorithm: Described elsewhere

\subsection*{Input parameters}
Parameters in \textbf{boldface} are required; the others are optional.

\begin{longtable}{p{0.22\textwidth}p{0.12\textwidth}p{0.46\textwidth}p{0.14\textwidth}}
\toprule
\textbf{Name} & \textbf{Unit} & \textbf{Description} & \textbf{Default} \\
\midrule
\endhead
density & g/cm$^{3}$ & Nominal density of material. & 0 \\
init & str & Name of Union inititaliser component & "init" \\
\bottomrule
\end{longtable}

\subsection*{Links}
\begin{itemize}
  \item Component source code found in file \texttt{KN\_xrl\_process.comp}.
  \item The test/example instrument \htmladdnormallink{Test\_Phonon.instr}{../examples/Test\_Phonon.instr}.
\end{itemize}
\IfFileExists{union/KN_xrl_process_static.tex}{\input{union/KN_xrl_process_static.tex}}{}
\input{union/Powder_process}
\section{The \texttt{Template\_process} McXtrace Component}
Template for a new contributor to create their own physical process.

\subsection*{Identification}
\begin{itemize}
  \item \textbf{Author:} Mads Bertelsen and Erik B Knudsen
  \item \textbf{Origin:} ESS DMSC \& DTU Physics
  \item \textbf{Date:} 20.08.15
\end{itemize}

\subsection*{Description}
This is a template for a new contributor to create their own physical process. The comments in this file are meant to teach the user about creating their own process file, rather than explaining this one. For comments on how this code works, look in the Incoherent\_process.comp.

Part of the Union components, a set of components that work together and thus sperates geometry and physics within McXtrace. The use of this component requires other components to be used.

1) One specifies a number of processes using process components like this one 2) These are gathered into material definitions using Union\_make\_material 3) Geometries are placed using Union\_box / Union\_cylinder, assigned a material 4) A Union\_master component placed after all of the above

Only in step 4 will any simulation happen, and per default all geometries defined before the master, but after the previous will be simulated here.

There is a dedicated manual available for the Union\_components

Algorithm: Described elsewhere

\subsection*{Input parameters}
Parameters in \textbf{boldface} are required; the others are optional.

\begin{longtable}{p{0.22\textwidth}p{0.12\textwidth}p{0.46\textwidth}p{0.14\textwidth}}
\toprule
\textbf{Name} & \textbf{Unit} & \textbf{Description} & \textbf{Default} \\
\midrule
\endhead
sigma & barns & Scattering cross section & 5.08 \\
packing\_factor & 1 & Material packing factor & 1 \\
unit\_cell\_volume & \AA{}$^{3}$ & Unit cell volume & 13.8 \\
interact\_fraction & 1 & How large a part of the scattering events should use this process 0-1 (sum of all processes in material = 1) & -1 \\
\bottomrule
\end{longtable}

\subsection*{Links}
\begin{itemize}
  \item Component source code found in file \texttt{Template\_process.comp}.
\end{itemize}
\IfFileExists{union/Template_process_static.tex}{\input{union/Template_process_static.tex}}{}
